% Vietnamese translation for OI Public Library
% Source-File: IOI中国国家候选队论文/2006/胡伟栋/胡伟栋（附）.ppt
\documentclass[11pt]{article}
\usepackage{oipl-vn}

\newcommand{\Oh}{\mathcal{O}}

\title{Bản trình chiếu phụ của Hu Weidong}
\author{Hu Weidong}
\date{2006}

\begin{document}
\maketitle

\origtitle{Hu Weidong appendix presentation}
\sourcepath{IOI China candidate team papers / 2006 / Hu Weidong / appendix presentation.ppt}

\section{Tình trạng nguồn}

Tệp phụ trích được nhiều mảnh hơn tệp chính: ngày 2006.1.19, một đường dẫn cục
bộ tới tài liệu Tsinsen, dãy độ phức tạp \(\Oh(n^5)\), \(\Oh(n^4)\),
\(\Oh(n^3)\), \(\Oh(n^2)\), \(\Oh(n\log^2 n)\), cùng vài công thức dạng
\texttt{Ans[0] = I} và \(F[4]=F[3]+F[2]\). Tuy nhiên phần thuyết minh quanh
các mảnh này không đọc được từ văn bản thuần.

\section{Ý nghĩa các mốc}

Dãy độ phức tạp cho thấy slide phụ nhiều khả năng minh họa quá trình cải tiến
thuật toán qua nhiều mức, từ nghiệm rất chậm tới nghiệm gần tuyến tính nhân
logarit. Công thức truy hồi với \(F[4]\) gợi ý một ví dụ đếm hoặc phân rã
kiểu Fibonacci, nơi cùng một giá trị được tách thành các giá trị nhỏ hơn.

\section{Các công thức truy hồi}

Phần cuối tệp có các công thức truy hồi đọc được:
\[
  F[4]=F[3]+F[2]=F[2]+F[1]+F[1]+F[1]
  =F[1]+F[1]+F[1]+F[1]+F[1]=5,
\]
cùng các dạng tổng quát
\[
  F[n]=2F[n-1],\qquad
  F[n]=F[n-1]+F[n-2],\qquad
  F[n]=\frac{F[n+1]+F[n-1]}2.
\]
Các công thức này cho thấy slide phụ so sánh nhiều cách xây dựng hoặc suy
diễn cùng một dãy, có thể nhằm minh họa tác động của mô hình truy hồi lên độ
phức tạp.

\section{Cách xử lý}

Do nội dung còn lại nằm trong ảnh, hiệu ứng hoặc đối tượng nhúng, bản dịch
không biến các mảnh rời thành một bài thuật toán hoàn chỉnh. Nó ghi lại các
tín hiệu kỹ thuật có thể xác nhận: dãy độ phức tạp, truy hồi Fibonacci, một
truy hồi nhân đôi, truy hồi trung bình hai phía, và ký hiệu mảng \texttt{Ans}.
Đây là mức phục dựng thận trọng nhất khi chưa có bản hiển thị từng trang.

\section{Khi có nguồn hiển thị}

Nếu có bản chuyển đổi sang ảnh, cần đối chiếu từng trang với các mốc độ phức
tạp trên để phục dựng đường dây cải tiến thuật toán. Khi đó phần này có thể
được thay bằng ghi chú chi tiết hơn về bài toán và các bước tối ưu.

\end{document}
